
\bibitem{src-292} J. Moser, Finitely many mass points on the line under the influence of an exponential potential -- an integrable system, in \emph{Dynamical Systems, Theory and Applications}, 467--497 Springer, Berlin, 1975.
\bibitem{src-293} D. Mumford, J. Fogarty, and F. Kirwan, \emph{Geometric Invariant Theory}, Springer, Berlin, 1994.
\bibitem{src-294} I. Mundet i Riera, Yang--Mills theory for symplectic fibrations, PhD thesis, Madrid, April 1999.
\bibitem{src-295} I. Mundet i Riera, Hamiltonian Gromov--Witten invariants, \emph{Topology} 42 (2003), 525--553.
\bibitem{src-296} V. Munoz, On the Gromov--Witten invariants of the moduli of bundles on a surface, AG/9910105, \emph{Forum Math.} 14 (2002), 759--774.
\bibitem{src-297} I. Naraki, Cross ratio variety as a moduli space of cubic surfaces, \emph{Proc. London Math. Soc.} 45 (1982), 1--30.
\bibitem{src-298} S. Nemirovski, Lagrangian Klein bottles in $\mathbb R^{2n}$, \emph{Geom. Funct. Anal.} 19 (2009), no. 3, 902--909.
\bibitem{src-299} K.L. Nguyen, C. Woodward, F. Ziltener, Morphisms of CohFT algebras and quantization of the Kirwan map, Preprint 2009, revised 2011. arXiv:0903.4459
\bibitem{src-300} A. Nijenhuis and W. Woolf, Some integration problems in almost-complex and complex manifolds, \emph{Annals of Math.} 77 (1963), 424--489.
\bibitem{src-301} S.P. Novikov, Multivalued functions and functionals -- an analogue of the Morse theory, \emph{Soviet Math. Dokl.} 24 (1981), 222--225.
\bibitem{src-302} Y.-G. Oh, Removal of boundary singularities of pseudo-holomorphic curves with Lagrangian boundary conditions, \emph{Comm. Pure Appl. Math.} 45 (1992), 121--139.
\bibitem{src-303} Y.-G. Oh, Floer cohomology of Lagrangian intersections discs and pseudoholomorphic discs I, II. \emph{Comm. Pure Appl. Math.}, 46 (1993), 949--94\sourcefootnote{原书第一篇页码范围止页印作 ``94''，明显不完整；无法仅据本书确定正确止页，故照录。}, 995--1012. III. Arnold--Givental conjecture, in \cite{src-181}, 555--574.
\bibitem{src-304} Y.-G. Oh, Floer cohomology, spectral sequences and the Maslov class of Lagrangian embeddings, \emph{Int. Math. Res. Notes} 7 (1996), 305--346.
\bibitem{src-305} Y.-G. Oh, Chain level Floer theory and Hofer’s geometry on the Hamiltonian diffeomorphism group, \emph{Asian J. Math.} 6 (2003), 579--624.
\bibitem{src-306} Y.-G. Oh, Spectral invariants, analysis of the Floer moduli space and geometry of the Hamiltonian diffeomorphism group, \emph{Duke Math. J.} 130 (2005), no. 2, 199--295.
\bibitem{src-307} Y.-G. Oh, Spectral invariants and length minimizing property of Hamiltonian paths, \emph{Asian J. Math.} 9 (2005), no. 1, 1--18.
\bibitem{src-308} Y.-G. Oh, The group of Hamiltonian homeomorphisms and $C^0$ symplectic topology, \emph{J. Symplectic Geom.} 5 (2007), no.2, 167--219.
\bibitem{src-309} C. Okonek and A. Teleman, Gauge theoretic equivariant Gromov--Witten invariants and the full Seiberg--Witten invariants of ruled surfaces, \emph{Commun. Math. Phys.} 227 (2002), 551--585.
\bibitem{src-310} C. Okonek and A. Teleman, Gauge theoretical Gromov--Witten invariants and virtual fundamental classes, \emph{The Fano conference}, 591--623, Univ. Torino, Turin, (2004).
\bibitem{src-311} A. Okounkov and R. Pandharipande, Gromov--Witten theory, Hurwitz numbers, and completed cycles, \emph{Ann. of Math.} (2) 163 (2006), no. 2, 517--560.
\bibitem{src-312} A. Okounkov and R. Pandharipande, The equivariant Gromov--Witten theory of $\mathbb P^1$, \emph{Ann. of Math.} (2) 163 (2006), no. 2, 561--605.
\bibitem{src-313} K. Ono, The Arnold conjecture for weakly monotone symplectic manifolds, \emph{Invent. Math.} 119 (1995), 519--537.
\bibitem{src-314} Y. Ostrover, Calabi quasi-morphisms for some non-monotone symplectic manifolds, \emph{Alg. Geom. Top.} 6 (2006), 405--434.
\bibitem{src-315} Y. Ostrover, A comparison of Hofer’s metrics on Hamiltonian diffeomorphisms and Lagrangian submanifolds, \emph{Commun. Contemp. Math} 5 (2003), no. 5, 803--811.
\bibitem{src-316} Y. Ostrover and I. Tyomkin, On the quantum homology algebra of toric Fano manifolds, \emph{Selecta Math.} (N.S.) 15 (2009), no. 1, 131--149.
\bibitem{src-317} Y. Ostrover and R. Wagner, On the extremality of Hofer’s metric on the group of Hamiltonian diffeomorphisms, \emph{Int. Math. Res. Not.} 2005, no 35, 2123--2141.
\bibitem{src-318} P. Ozsváth and Z. Szabó, Holomorphic discs and three-manifold invariants for closed 3-manifolds, \emph{Ann. Math.} 159 (2004), 1027--1158.
\bibitem{src-319} A. Ott. Removal of singularities and Gromov compactness for symplectic vortices, arXiv:0912.2500, to appear in \emph{J. Symplectic Geom.}
\bibitem{src-320} P. Pansu, Pseudo-holomorphic curves in symplectic manifolds, in \cite{src-25}, 233--250.
\bibitem{src-321} T.H. Parker and J.G. Wolfson, Pseudoholomorphic maps and bubble trees, \emph{J. Geom. Anal.} 3 (1993), 63--98.
\bibitem{src-322} M. Pinsonnault, Symplectomorphism groups and embeddings of balls into rational ruled 4-manifolds, \emph{Compos. Math.} 144 (2008), no. 3, 787--810.
\bibitem{src-323} S. Piunikhin, D. Salamon, M. Schwartz Symplectic Floer-Donaldson theory and quantum cohomology, in \emph{Contact and Symplectic Geometry}, edited by C.B. Thomas, Publications of the Newton Institute, Cambridge University Press 1996, 171--200.
\bibitem{src-324} L. Polterovich, Monotone Lagrangian submanifolds of linear spaces and the Maslov class in cotangent bundles, \emph{Math. Zeit.} 207 (1991), 217--222.
\bibitem{src-325} L. Polterovich, Gromov’s K-area and symplectic rigidity, \emph{Geom. Funct. Anal.} 6 (1996), 726--739.
\bibitem{src-326} L. Polterovich, Symplectic aspects of the first eigenvalue, \emph{J. Reine Angew. Math.} 502 (1998), 1--17.
\bibitem{src-327} L. Polterovich, Precise measurements in symplectic topology, in \emph{European Congress of Mathematics}, Vol. II (Budapest, 1996), 159--166, Progr.Math. 169, Birkhäuser, Basel, 1998.
\bibitem{src-328} L. Polterovich, Geometry on the group of Hamiltonian diffeomorphisms, Proceedings of the International Congress of Mathematicians, (Berlin 1998), \emph{Doc. Math.}, Extra Vol. II, 401--410.
\bibitem{src-329} L. Polterovich, Growth of maps, distortion in groups and symplectic geometry, \emph{Invent. Math.} 150 (2002), 655--686.
\bibitem{src-330} L. Polterovich, \emph{The Geometry of the Group of Symplectic Diffeomorphisms}, Lectures in Math, ETH, Birkhauser, 2001.
\bibitem{src-331} L. Polterovich, Quantum unsharpness and symplectic rigidity, arXiv:1110.5247
\bibitem{src-332} P. Py, Quasimorphisms and Calabi invariant, \emph{Ann. Sci. École Norm. Sup.} (4) 39 (2006), 177--194.
\bibitem{src-333} J.W. Robbin, Y. Ruan, and D.A. Salamon, The moduli space of regular stable maps \emph{Math. Zeit.} 259 (2008), 525--574.
\bibitem{src-334} J.W. Robbin and D.A. Salamon, The Maslov index for paths, \emph{Topology} 32 (1993), 827--844.
\bibitem{src-335} J.W. Robbin and D.A. Salamon, The spectral flow and the Maslov index, \emph{Bull. L.M.S.} 27 (1995), 1--33.
\bibitem{src-336} J.W. Robbin and D.A. Salamon, Asymptotic behaviour of holomorphic strips, \emph{Annales de l’Institut Henri Poincaré -- Analyse Nonlineaire} 18 (2001), 573--612.
\bibitem{src-337} J.W. Robbin and D.A. Salamon, A construction of the Deligne-Mumford orbifold, \emph{J. Europ. Math. Soc.} 8 (2006), 611--699.
\bibitem{src-338} H.L. Royden, \emph{Real Analysis}, 3rd Edition, Macmillan, New York, 1988.
\bibitem{src-339} Y. Ruan, Symplectic topology and extremal rays, \emph{Geom. Funct. Anal.} 3 (1993), 395--430.
\bibitem{src-340} Y. Ruan, Symplectic Topology on Algebraic 3-folds, \emph{J. Diff. Geom.} 39 (1994), 215--227.
\bibitem{src-341} Y. Ruan, Symplectic topology and complex surfaces, in: \emph{Geometry and Analysis on Complex manifolds}, ed Mabuchi, Noguchi, Ochiai, World Scientific (1994), 171--197.
\bibitem{src-342} Y. Ruan, Topological sigma model and Donaldson type invariants in Gromov theory, \emph{Duke Math J.} 83 (1996), 461--500.
\bibitem{src-343} Y. Ruan, Virtual neighborhoods and pseudoholomorphic curves, in \cite{src-385}, 101--116.
\bibitem{src-344} Y. Ruan, Surgery, quantum cohomology, and birational geometry, in \cite{src-99}, 183--199.
\bibitem{src-345} Y. Ruan and G. Tian, A mathematical theory of quantum cohomology, \emph{Math. Res. Letters} 1 (1994), 269--278.
\bibitem{src-346} Y. Ruan and G. Tian, A mathematical theory of quantum cohomology, \emph{J. Diff. Geom.} 42 (1995), 259--367.
\bibitem{src-347} Y. Ruan and G. Tian, Bott--type symplectic Floer cohomology and its multiplicative structures, \emph{Math. Res. Letters}, 2 (1995), 203--219.
\bibitem{src-348} Y. Ruan and G. Tian, Higher genus symplectic invariants and sigma model coupled with gravity, \emph{Invent. Math.} 130 (1997), 455--516.
\bibitem{src-349} J. Sacks and K. Uhlenbeck, The existence of minimal immersions of 2-spheres, \emph{Annals of Math.} 113 (1981), 1--24.
\bibitem{src-350} V. Sadov, On equivalence of Floer’s and quantum cohomology, hep-th/9310153, \emph{Commun. Math. Phys.} 173 (1995), 77--100.
\bibitem{src-351} D.A. Salamon, Morse theory, the Conley index and Floer homology, \emph{Bull. London.M.S.} 22 (1990), 113--140.
\bibitem{src-352} D.A. Salamon, Quantum products for mapping tori and the Atiyah-Floer conjecture, \emph{Amer. Math. Soc. Transl.} 196 (1999), 199--235.
\bibitem{src-353} D.A. Salamon, Lectures on Floer homology, in \cite{src-103}, 145--229.
\bibitem{src-354} D.A. Salamon, A survey of symmetric functions, Grassmannians, and representations of the unitary group. Unpublished Notes, University of Warwick, 1996. \url{http://www.math.ethz.ch/salamon/publications.html}
\bibitem{src-355} D.A. Salamon and K. Wehrheim, Instanton Floer homology with Lagrangian boundary conditions. \emph{Geometry \& Topology} 12 (2008), 747--918.
\bibitem{src-356} D.A. Salamon and E. Zehnder, Morse theory for periodic solutions of Hamiltonian systems and the Maslov index, \emph{Comm. Pure Appl. Math.} 45 (1992), 1303--1360.
\bibitem{src-357} M. Schwarz, \emph{Morse Homology}, Birkhäuser, 1993.
\bibitem{src-358} M. Schwarz, Cohomology operations from $S^1$-cobordisms in Floer theory, PhD thesis, ETH-Zürich, 1995.
\bibitem{src-359} M. Schwarz, Equivalences for Morse homology, in \emph{Geometric Topology in Dynamics}, ed. Barge and Kuperberg, \emph{Contemporary Mathematics} 246, Amer. Math. Soc. 1999.
\bibitem{src-360} M. Schwarz, On the action spectrum for closed symplectically aspherical manifolds, \emph{Pacific J. Math} 193 (2000), 419--461.
\bibitem{src-361} P. Seidel, Floer homology of a Dehn twist, \emph{Math. Res. Letters} 3 (1997), 829--834.
\bibitem{src-362} P. Seidel, \emph{Floer homology and the symplectic isotopy problem}, PhD thesis, Oxford, 1997.
\bibitem{src-363} P. Seidel, $\pi_1$ of symplectic automorphism groups and invertibles in quantum homology rings, \emph{Geom. Funct. Anal.} 7 (1997), 1046--1095.
\bibitem{src-364} P. Seidel, On the length of Hamiltonian loops, unpublished preprint (1997).
\bibitem{src-365} P. Seidel, On the group of symplectic automorphisms of $\mathbb CP^m\times\mathbb CP^n$, in \cite{src-99}, 237--250.
\bibitem{src-366} P. Seidel, Vanishing cycles and mutations, in: \emph{Third European Congress of Mathematics} (Barcelona 2000), Progress in Mathematics 202 Birkhäuser (2001), 65--86.
\bibitem{src-367} P. Seidel, More on vanishing cycles and mutations, in: \emph{Symplectic Geometry and Mirror Symmetry} (Seoul 2000), World Scientific (2001), 249--265.
\bibitem{src-368} P. Seidel, Symplectic automorphisms of $T^*S^2$, arXiv:DG/9803084.
\bibitem{src-369} P. Seidel, Symplectic Floer homology and the mapping class group, \emph{Pac. J. Math.} 206 (2002), 219--229.
\bibitem{src-370} P. Seidel, A long exact sequence for symplectic Floer homology, \emph{Topology} 42 (2003), 1003--1064.
\bibitem{src-371} P. Seidel, \emph{Fukaya Categories and Picard--Lefschetz theory}, Zurich Lectures in Advanced Mathematics, European Math. Soc. (EMS), Zurich, 2008.
\bibitem{src-372} P. Seidel, Lectures on four-dimensional Dehn twists, \emph{Symplectic four manifolds and algebraic surfaces}, 231--267, Lecture Notes in Math, 1938, Springer, Berlin, 2008.
\bibitem{src-373} P. Seidel, Abstract analogs of flux as symplectic invariants, arXiv:1108.0394.
\bibitem{src-374} E. Shelukhin, The action homomorphism, quasimorphisms, and moment maps on the space of almost complex structures, arXiv:1105.5814.
\bibitem{src-375} V. V. Shevchishin, Pseudoholomorphic curves and the symplectic isotopy problem, Habilitation thesis, Ruhr-Univ, Bochum, arXiv:0010262.
\bibitem{src-376} B. Siebert, Symplectic Gromov--Witten invariants, in: \emph{New Trends in Algebraic Geometry}, ed. Catanese, Peters, Reid, L.M.S. Lect. Notes 264, Camb. Univ. Press 1999, 375--424.
\bibitem{src-377} B. Siebert and G. Tian, Recursive relations for the cohomology ring of moduli spaces of stable bundles, \emph{Turkish Journal of Mathematics} 19 (1995), 131--144.
\bibitem{src-378} B. Siebert and G. Tian, On quantum cohomology rings of Fano manifolds and a formula of Vafa and Intriligator, \emph{Asian J. Math.} 1 (1997), 679--695.
\bibitem{src-379} J.-C. Sikorav, Singularities of $J$-holomorphic curves, \emph{Math. Zeitschrift} 226 (1997), 359--373.
\bibitem{src-380} J.-C. Sikorav, Dual elliptic planes, \emph{Actes des Journées Math. à la mémoire de Jean Leray}, 185--207, Sém. Congr., 9, Soc. Math. France, Paris, 2004.
\bibitem{src-381} S. Smale, An infinite dimensional version of Sard’s theorem, \emph{Am. J. Math.} 87 (1973), 213--221.
\bibitem{src-382} I. Smith, On moduli spaces of symplectic forms, \emph{Math. Res. Letters} 7 (2000), 779--788.
\bibitem{src-383} I. Smith, Torus fibrations on symplectic 4-manifolds, \emph{Turkish J. Math.} 25 (1) (2001), 69--95.
\bibitem{src-384} H. Spielberg, Counting Generic Genus 0 curves on Hirzebruch surfaces, \emph{Proc. Amer. Math. Soc.} 130 (2001), 1257--1264.
\bibitem{src-385} R. Stern, ed., \emph{Topics in Symplectic 4-manifolds} (Irvine CA 1996), Internat. Press Lecture Series, Vol 1, Internat. Press, Cambridge, MA, 1998.
\bibitem{src-386} J. Streets and G. Tian, Symplectic Curvature Flow, arXiv:1012.2104.
\bibitem{src-387} Z. Szabo, Simply connected irreducible 4-manifolds with no symplectic structure, \emph{Invent. Math.} 132 (1998), 457--466.
\bibitem{src-388} D. Tamarkin, Microlocal condition for non-displaceability, arXiv:0809.1584
\bibitem{src-389} C.H. Taubes, Counting pseudoholomorphic submanifolds in dimension 4, \emph{J. Diff. Geom.} 44 (1996), 818--893. (republished in \cite{src-393})
\bibitem{src-390} C.H. Taubes, The Seiberg-Witten invariants and symplectic forms, \emph{Math. Res. Lett.} 1 (1994), 809--822.
\bibitem{src-391} C.H. Taubes, More constraints on symplectic forms from Seiberg-Witten invariants, \emph{Math. Res. Lett.} 2 (1995), 9--14.
\bibitem{src-392} C.H. Taubes, The Seiberg-Witten and the Gromov invariants, \emph{Math. Res. Lett.} 2 (1995), 221--238.
\bibitem{src-393} C.H. Taubes, \emph{Seiberg--Witten and Gromov Invariants for Symplectic 4- Manifolds}, First Internat. Press Lecture Series vol 2, Internat. Press. Cambridge MA, 2000.
\bibitem{src-394} C.H. Taubes, The Seiberg--Witten equations and the Weinstein conjecture, \emph{Geom. Topol.} 11 (2007), 2117--2202.
\bibitem{src-395} R. Thom, Quelques propriétés globales des variétés différentiables, \emph{Comment. Math. Helv.} 29 (1954), 17--85.
\bibitem{src-396} Zhiyu Tian, Some examples of symplectic rationally connected 4-folds, arXiv:1112.1369.
\bibitem{src-397} V. Tosatti and B. Weinkove, The Calabi--Yau equation on the Kodaira--Thurston manifold, \emph{J. Inst. Math. Jussieu} 10 (2011), 437--447.
\bibitem{src-398} M. Usher, Spectral numbers in Floer theories, \emph{Compos. Math.} 144 (2008), 1581--1592.
\bibitem{src-399} M. Usher, Deformed Hamiltonian Floer theory, capacity estimates and Calabi quasimorphisms, \emph{Geom. Topol.} 15 (2011), 1313--1417.
\bibitem{src-400} M. Usher, Hofer’s metrics and boundary depth, arXiv:1107.4599
\bibitem{src-401} C. Vafa, Topological mirrors and quantum rings, in: \emph{Essays on Mirror Manifolds}, ed. S.-T. Yau, International Press, Hong Kong 1992.
\bibitem{src-402} C. Viterbo, A proof of the Weinstein conjecture in $\mathbb R^{2n}$, \emph{Annales de l’ Institut Henri Poincaré -- Analyse Nonlinéaire} 4 (1987), 337--357.
\bibitem{src-403} C. Viterbo, A new obstruction to embedding Lagrangian tori, \emph{Invent. Math.} 100 (1990), 301--320.
\bibitem{src-404} C. Viterbo, Symplectic topology as the geometry of generating functions. \emph{Math. Annalen}, 292 (1992), 685--710.
\bibitem{src-405} C. Viterbo, The cup-product on the Thom-Smale-Witten complex, and Floer cohomology, in \cite{src-181}, 609--626.
\bibitem{src-406} C. Viterbo, Functors and computations in Floer homology with applications, part I, \emph{Geom. Funct. Anal.} 9 (1999), 985--1033.
\bibitem{src-407} C. Viterbo, On the uniqueness of generating Hamiltonian for continuous limits of Hamiltonian flows, \emph{Int. Math. Res. Not} 2006, Art ID 34028, erratum Art ID 38784
\bibitem{src-408} C. Voisin, A mathematical proof of a formula of Aspinwall and Morrison, \emph{Compositio Math.}, 104, (1996), 135--151.
\bibitem{src-409} C. Voisin, \emph{Mirror Symmetry}, SMF/AMS texts and Monographs, Vol 1, Amer. Math. Soc., Providence, RI, 1996.
\bibitem{src-410} C. Voisin, Rationally connected 3-folds and symplectic geometry, \emph{Astérisque} 322 (2008), 1--21.
\bibitem{src-411} K. Wehrheim, Anti-self-dual instantons with Lagrangian boundary conditions, PhD thesis, ETH-Zürich, July 2002.
\bibitem{src-412} K. Wehrheim, Banach space valued Cauchy--Riemann equations with totally real boundary conditions, \emph{Commun. Contemp. Math.} 6 (2004), no. 4, 601--635.
\bibitem{src-413} K. Wehrheim, Anti-self-dual instantons with Lagrangian boundary conditions: elliptic theory, \emph{Comm. Math. Phys.} 254 (2005), no. 1, 45--89.
\bibitem{src-414} K. Wehrheim and C. Woodward, Functoriality for Lagrangian correspondences in Floer theory, \emph{Quantum Topol.} 1 (2010) no. 2, 129--170.
\bibitem{src-415} K. Wehrheim and C. Woodward, Floer Cohomology and Geometric Composition of Lagrangian Correspondences, arXiv:0905.1368
\bibitem{src-416} K. Wehrheim and C. Woodward, Pseudoholomorphic Quilts, arXiv:0905.1369
\bibitem{src-417} K. Wehrheim and C. Woodward, Quilted Floer Cohomology, \emph{Geom. Topol.} 14 (2010), no. 2, 192--210.
\bibitem{src-418} C. Wendl, Strongly fillable contact manifolds and $J$-holomorphic foliations, \emph{Duke Math. J.} 151 (2010), no.3, 337--384.
\bibitem{src-419} P.M.H. Wilson, The Kähler cone on Calabi-Yau threefolds, \emph{Invent. Math.} 107 (1992), 561--584.
\bibitem{src-420} E. Witten, Supersymmetry and Morse theory, \emph{J. Diff. Geom.} 17 (1982), 661--692.
\bibitem{src-421} E. Witten, Two dimensional gravity and intersection theory on moduli space, \emph{Surveys in Diff. Geom.} 1 (1991), 243--310.
\bibitem{src-422} E. Witten, The Verlinde algebra and the cohomology of the Grassmannian, hep-th/9312104, in \emph{Geometry, Topology and Physics} Conf. Proc. Lecture Notes Geom. Topology VI, Internat. Press, Cambridge MA, 1995, 357--422.
\bibitem{src-423} R. Ye, Gromov’s compactness theorem for pseudoholomorphic curves, \emph{Trans. Amer. Math. Soc.} 342 (1994), 671--694.
\bibitem{src-424} J.-C. Yoccoz, Travaux de Herman sur les tores invariants, Sem. Bourbaki, 1991-2, no 754, \emph{Astérisque} 206 (1992), 311--344.
\bibitem{src-425} K. Zehmisch, The annulus property of simple holomorphic discs, \emph{J. of Symplectic Geom.} (2012).
\bibitem{src-426} A. Zinger, Enumerative vs. Symplectic Invariants and Obstruction Bundles, \emph{J. Symplectic Geom.} 2 (2004), no. 4, 445--543.
\bibitem{src-427} A. Zinger, Enumeration of genus 2 curves with a fixed complex structure in $\mathbb P^2$ and $\mathbb P^3$, \emph{J. Differential Geom.} 65 (2003) no. 3, 341--467.
\bibitem{src-428} A. Zinger, Counting rational curves of arbitrary shape in projective spaces, \emph{Geom. Topol.} 9 (2005), no. 1, 35--81.
\bibitem{src-429} A. Zinger, Pseudocycles and integral homology, \emph{Trans. Amer. Math. Soc.} 360 (2008), no. 5, 2741--2765.
\bibitem{src-430} F. Ziltener, Symplectic vortices on the complex plane and quantum cohomology, Ph.D. thesis, ETH Zürich, 2006.
\bibitem{src-431} F. Ziltener, The Invariant Symplectic Action and Decay for Vortices, \emph{J. Symplectic Geom.} 7 (2009) no. 3, 357--376.
\bibitem{src-432} F. Ziltener, A Quantum Kirwan Map: Bubbling and Fredholm Theory, arXiv:0905.4047 and arXiv:1106.1729, to appear in \emph{Mem. Amer. Math. Soc.}
